% !PRINTED 436
\weberpage{436}
Die sämtlichen Wurzeln von (3) sind aber nach § 69:
\weq{5}{j(p\omega), \quad j\left(\frac{\omega + \lambda}{p}\right),}
wenn $\lambda$ ein Restsystem nach dem Modul $p$ durchläuft. Die erste
von diesen gehört nicht zu den Wurzeln von (4), denn $\omega' = p\omega$
genügt der Gleichung:
\[
a\omega'^2 + p b \omega' + p^{\varepsilon+2} c = 0,
\]
die primitiv und von der Diskriminante $-mp^2$ ist. Wenn wir aber
$p\omega' = \omega + \lambda$ setzen, so erhalten wir für $\omega'$ die Gleichung:
\weq{6}{a p^2 \omega'^2 + p(b - 2a\lambda)\omega' + (a\lambda^2 - b\lambda + c p^{\varepsilon}) = 0,}
und diese Gleichungen sind ebenfalls von der Diskriminante
$-mp^2$. Es sind aber zwei darunter, die imprimitiv vom Teiler $p$
sind, nämlich die den Werten
\weq{7}{\lambda \equiv 0, \quad a\lambda - b \equiv 0 \pmod{p}}
entsprechen. Die erste ist die Wurzel der Form $(a p, b, c p^{\varepsilon-1})$, die
zweite die der Form
\[
(a p, \; b - 2a\lambda, \; C),
\]
und da $b - 2a\lambda \equiv b \pmod{2a}$ und [nach (7)] $\equiv -b \pmod{2p}$
ist, so ist diese Form komponiert aus $(a, b, c p^{\varepsilon})$ und $(p, -b, a c p^{\varepsilon-1})$,
die in die zu $l$ reziproke Klasse $l^{-1}$ gehört. Daraus ergibt sich:

I. \gesperrt{Die Gleichungen (3) und (4) haben nur zwei Wurzeln
 miteinander gemein, nämlich die Klasseninvarianten
$(lk)$ und $(l^{-1}k)$, und wenn die Klasse $l$
zweiseitig ist, so haben sie nur eine Wurzel
gemein.}

Setzt man also in (3) $u = (k)$, so ist der größte gemeinschaftliche
 Teiler von (3) und (4) vom zweiten, und wenn $l$
zweiseitig ist, vom ersten Grade. Im ersteren Falle sind $(lk)$ und
$(l^{-1}k)$ die Wurzeln einer quadratischen Gleichung, deren Koeffizienten
 rational von $u$ abhängen, und deren Form im übrigen
nur von der Klasse $l$, nicht von $k$ abhängig ist. Dem hiermit
bewiesenen können wir den Ausdruck geben:

II. \gesperrt{Es ist}
\weq{8}{(lk) + (l^{-1}k) = f_l(k),}
\gesperrt{oder, wenn $l$ zweiseitig ist,}
\weq{9}{(lk) = f_l(k),}
\gesperrt{worin $f_l(k)$ eine rationale Funktion von $(k)$ bedeutet,
deren rationale Zahlenkoeffizienten nur
von der Klasse $l$, nicht von der Klasse $k$ abhängen.}
